% Manuscript-ready Methods text derived from
% EQUIVALENCE_THEOREM_FITTING_REPORT.tex.
%
% Required packages in the parent manuscript:
%   \usepackage{amsmath,amssymb,bm}
%   \usepackage[round,authoryear]{natbib}  % or adapt citation commands
%
% Required bibliography entries:
%   irvine1964, partain2000, stephens2000, heidinger2000,
%   funk2003, min2004, scholl2006, bennartz2006

% -----------------------------------------------------------------------------
% MAIN-TEXT METHODS
% -----------------------------------------------------------------------------

\subsection{Spectrum-internal photon-path features}
\label{sec:methods_photon_path}

We derived spectral features that summarize the distribution of atmospheric
photon paths represented in each OCO-2 absorption band.  The derivation follows
the radiative-transfer equivalence theorem
\citep{irvine1964,partain2000,stephens2000}.  For a fixed illumination and
viewing geometry, let $P_0(\Gamma)$ denote the ensemble of photon trajectories
$\Gamma$ reaching the detector in a conservative reference atmosphere with the
same surface boundary and continuum scattering properties as the absorbing
atmosphere.  Across each narrow OCO-2 band, we assume that these scattering
properties vary slowly relative to molecular line absorption.  Wavelength then
changes the contribution of a trajectory primarily through its Beer--Lambert
survival probability,
\begin{equation}
  w_\lambda(\Gamma)
  = \exp[-A_\lambda(\Gamma)],
  \qquad
  A_\lambda(\Gamma)
  = \int_\Gamma \alpha_\lambda(\bm r)\,\mathrm{d}s,
  \label{eq:main_path_weight}
\end{equation}
where $\alpha_\lambda$ is the local molecular absorption coefficient.  The
normalized radiance is therefore the conservative-ensemble mean
\begin{equation}
  \mathcal T_\lambda
  = \int P_0(\Gamma)
      \exp[-A_\lambda(\Gamma)]\,\mathrm{d}\Gamma.
  \label{eq:main_equivalence}
\end{equation}
This construction separates the multiple-scattering calculation from molecular
absorption and remains applicable to three-dimensional cloud fields because the
cloud geometry is encoded in $P_0(\Gamma)$ \citep{stephens2000}.

To obtain a tractable scalar representation, we approximate the absorber amount
along each trajectory as
\begin{equation}
  A_\lambda(\Gamma)
  \simeq \tau_\lambda\ell'(\Gamma),
  \label{eq:main_scalar_path}
\end{equation}
where $\tau_\lambda$ is the modeled molecular optical depth along the direct
geometrical slant path and $\ell'$ is an effective absorber-path enhancement
relative to that reference.  The conservative ensemble induces a normalized
distribution $p_0(\ell')$, so Eq.~\eqref{eq:main_equivalence} reduces to
\begin{equation}
  \mathcal T(\tau_\lambda)
  = \int_0^\infty p_0(\ell')
      \exp(-\tau_\lambda\ell')\,\mathrm{d}\ell'.
  \label{eq:main_laplace}
\end{equation}
Thus, the resolved absorption spectrum samples the Laplace transform of one
wavelength-independent conservative path distribution within each band.  This
does not imply that weak and strong channels detect identical photon
populations: absorption reweights the detected paths as
$p_{\mathrm{det}}(\ell'\mid\tau)\propto
p_0(\ell')\exp(-\tau\ell')$, preferentially suppressing long paths at large
$\tau$.

We normalized the measured radiance as
\begin{equation}
  T_\lambda
  = \frac{\pi I_\lambda}{F_{0,\lambda}\mu_0}
  = \gamma\,\mathcal T(\tau_\lambda),
  \label{eq:main_transmittance}
\end{equation}
where $F_{0,\lambda}$ is the top-of-atmosphere solar irradiance, $\mu_0$ is the
cosine of the solar zenith angle, and $\gamma$ is an effective continuum
reflectance.  In practice, $\gamma$ can include surface bidirectional
reflectance, polarization, and surface--atmosphere coupling and is not
interpreted as an exact Lambertian albedo.

Let $C_n$ denote the ordinary cumulants of $p_0(\ell')$.  Expanding the log
Laplace transform about zero absorption gives
\begin{equation}
  \ln T(\tau)
  = \ln\gamma
    + \sum_{n=1}^{\infty}
      \frac{(-1)^n C_n}{n!}\tau^n.
  \label{eq:main_cumulant}
\end{equation}
For each sounding and band, we fitted the truncated form
\begin{equation}
  \ln T(\tau)
  \simeq b
    + \sum_{n=1}^{N}
      \frac{(-1)^n k_n}{n}\tau^n,
  \qquad
  k_n=\frac{C_n}{(n-1)!},
  \label{eq:main_fit}
\end{equation}
where $b$ is the fitted continuum intercept.  Under the equivalence-theorem and
scalar-path assumptions, $k_1=\mathrm{E}_{p_0}[\ell']$ is the mean absorber-path
enhancement and $k_2=\mathrm{var}_{p_0}(\ell')$ is its variance; coefficients
$k_{n\geq3}$ are scaled higher cumulants.

The model is linear in $(b,k_1,\ldots,k_N)$.  We solved it by singular-value
least squares and used bounded-variable least squares only when the
unconstrained solution violated $k_1,k_2\geq0$.  We excluded the two edge
channels and channels with $T_\lambda>1$.  The truncation orders were 7, 3, and
7 for the O$_2$ A, weak-CO$_2$, and strong-CO$_2$ bands, respectively; the lower
weak-CO$_2$ order reflects its smaller optical-depth range and the resulting
poor identifiability of high-order terms.  Fits were produced for both raw and
Savitzky--Golay-smoothed log transmittance; the raw coefficients were used in
the production model, while the smoothed coefficients were retained for a
robustness comparison.

We interpret these quantities as \emph{band-effective absorber-path cumulants}.
The scalar approximation in Eq.~\eqref{eq:main_scalar_path} can fail when
pressure- and temperature-dependent line absorption weights atmospheric layers
differently, or when surface, aerosol, or polarization properties vary within a
band.  A layer-resolved formulation and the distinction from a gamma path model
are given in Appendix~\ref{app:photon_path_derivation}.


% -----------------------------------------------------------------------------
% APPENDIX METHODS
% -----------------------------------------------------------------------------

\section{Photon-path derivation, coefficient convention, and limitations}
\label{app:photon_path_derivation}

\subsection{Equivalence theorem and conservative path ensemble}
\label{app:equivalence}

The equivalence theorem separates the statistics of photon transport from the
spectral absorption applied along those paths \citep{irvine1964}.  In its most
general form relevant here, the normalized radiance is
\begin{equation}
  \mathcal T_\lambda
  = \mathrm{E}_{P_0}
    \left\{
      \exp[-A_\lambda(\Gamma)]
    \right\},
  \qquad
  A_\lambda(\Gamma)
  = \int_\Gamma \alpha_\lambda(\bm r)\,\mathrm{d}s.
  \label{eq:app_general_equivalence}
\end{equation}
The theorem requires absorption to modify photon weights without changing the
geometry, scattering phase function, or boundary conditions that generate
$P_0(\Gamma)$.  It does not require a plane-parallel medium.

\citet{stephens2000} used this construction for molecular-line absorption in
heterogeneous clouds.  A continuum Monte Carlo calculation supplied a photon
optical-path distribution and continuum radiance at each spatial location; an
entire high-resolution O$_2$ A-band spectrum was then reconstructed by applying
the wavelength-dependent transmission function to those stored path
statistics.  Direct multiple-scattering and path-distribution reconstructions
agreed in their test, while changes between plane-parallel and heterogeneous
clouds appeared through changes in the path distribution itself.  Their O$_2$
A-band application subsequently demonstrated sensitivity to cloud optical
depth, phase function, cloud pressure and pressure thickness, and surface
albedo \citep{heidinger2000}.

\citet{partain2000} generalized the computational use of the theorem by
retaining geometrical path and scattering statistics that allow spectral gas,
particle, and surface absorption to be applied after the multiple-scattering
calculation.  Their treatment accommodates multiple homogeneous layers and
vertical gas profiles.  The full joint distribution is memory-intensive,
however, and reduced-PDF approximations discard trajectory information.  The
scalar distribution fitted here is deliberately simpler than that generalized
bookkeeping: it retains one effective absorber-path coordinate per band rather
than the complete layer-resolved trajectory history.

\subsection{Scalar path reduction and detected-path selection}
\label{app:scalar_path}

Let $\bar L_z$ be the direct geometrical path through atmospheric layer $z$, and
let $L_z(\Gamma)$ be the corresponding path accumulated by trajectory
$\Gamma$.  The exact layered absorber amount is
\begin{equation}
  A_\lambda(\Gamma)
  = \sum_z \alpha_{\lambda z}L_z(\Gamma).
  \label{eq:app_layer_absorption}
\end{equation}
The one-dimensional model assumes approximately proportional layer paths,
\begin{equation}
  L_z(\Gamma)\simeq\ell'(\Gamma)\bar L_z,
  \label{eq:app_collinear_paths}
\end{equation}
which yields
\begin{equation}
  A_\lambda(\Gamma)
  \simeq \ell'(\Gamma)
          \sum_z\alpha_{\lambda z}\bar L_z
  = \tau_\lambda\ell'(\Gamma).
  \label{eq:app_scalar_absorption}
\end{equation}
The conservative trajectory ensemble then induces
\begin{equation}
  p_0(\ell')
  = \int P_0(\Gamma)
    \delta[\ell'-\ell'(\Gamma)]\,\mathrm{d}\Gamma,
  \label{eq:app_pdf}
\end{equation}
and Eq.~\eqref{eq:app_general_equivalence} becomes the one-dimensional Laplace
transform
\begin{equation}
  \mathcal T(\tau)
  = \int_0^\infty p_0(\ell')e^{-\tau\ell'}\,\mathrm{d}\ell'.
  \label{eq:app_laplace}
\end{equation}

The reference $\ell'=1$ represents the modeled direct surface-slant absorber
path.  Multiple scattering can give $\ell'>1$, whereas reflection above the
surface can truncate the path and give $0\leq\ell'<1$.  Accordingly, the fit
requires nonnegative, not necessarily greater-than-unity, mean path
enhancement.

The distribution among photons detected at a finite optical depth is
\begin{equation}
  p_{\mathrm{det}}(\ell'\mid\tau)
  = \frac{p_0(\ell')e^{-\tau\ell'}}
         {\mathcal T(\tau)}.
  \label{eq:app_detected_pdf}
\end{equation}
It follows that
\begin{align}
  -\frac{\mathrm{d}\ln\mathcal T}{\mathrm{d}\tau}
  &= \mathrm{E}_{\mathrm{det},\tau}[\ell'],
  \label{eq:app_detected_mean}\\
  \frac{\mathrm{d}^2\ln\mathcal T}{\mathrm{d}\tau^2}
  &= \mathrm{var}_{\mathrm{det},\tau}(\ell').
  \label{eq:app_detected_variance}
\end{align}
Weak and strong channels therefore observe different absorption-selected path
populations even when they sample one common conservative $p_0(\ell')$.

High-resolution O$_2$ A-band measurements provide an observational basis for
this representation.  \citet{min2004} treated the resolved line spectrum as an
inverse-Laplace problem and emphasized the requirements imposed by spectral
resolution, slit-function characterization, stray light, and out-of-band
rejection.  Their O$_2$ and nearby water-vapor results also showed that
different absorber profiles yield different effective path moments despite
similar scattering wavelengths.  \citet{scholl2006} represented paths above,
within, and below cloud separately, used pressure- and temperature-dependent
O$_2$ absorption, and convolved the modeled spectrum with the measured slit
function.  Their analysis concentrated on the first two moments and generally
used a gamma path PDF.  These studies motivate the present moment-based fit but
also show that its scalar path is an absorber-weighted approximation, not a
complete reconstruction of geometrical trajectories.

\subsection{Cumulant convention used by the estimator}
\label{app:cumulants}

Let $C_n$ be the ordinary cumulants of $p_0(\ell')$.  Since
$\mathcal T(\tau)$ is the moment-generating function evaluated at $-\tau$,
\begin{equation}
  \ln\mathcal T(\tau)
  = \sum_{n=1}^{\infty}
    \frac{(-1)^nC_n}{n!}\tau^n.
  \label{eq:app_cumulant_series}
\end{equation}
Including the effective continuum factor gives
\begin{equation}
  \ln T(\tau)
  = \ln\gamma
    + \sum_{n=1}^{\infty}
      \frac{(-1)^nC_n}{n!}\tau^n.
  \label{eq:app_cumulant_reflectance}
\end{equation}
The estimator is parameterized as
\begin{equation}
  \ln T(\tau)
  \simeq b+sum_{n=1}^{N}
    \frac{(-1)^nk_n}{n}\tau^n,
  \label{eq:app_project_parameterization}
\end{equation}
and therefore
\begin{equation}
  b=\ln\gamma,
  \qquad
  k_n=\frac{C_n}{(n-1)!}.
  \label{eq:app_coefficient_mapping}
\end{equation}
Thus $k_1=C_1$ and $k_2=C_2$ are exactly the mean and variance, whereas
$k_3=C_3/2!$, $k_4=C_4/3!$, and higher $k_n$ are scaled cumulants.  This
factorial convention follows from the $1/n$ factors used in the design matrix
and should be retained when comparing the fitted coefficients with standard
probability notation.

The finite-order model uses the design matrix
\begin{equation}
  \bm X(\tau)
  = \left[1,-\tau,\frac{\tau^2}{2},-\frac{\tau^3}{3},\ldots,
           \frac{(-1)^N\tau^N}{N}\right].
  \label{eq:app_design_matrix}
\end{equation}
Only $k_1$ and $k_2$ are constrained nonnegative.  This is necessary for their
mean and variance interpretation, while higher cumulants of a general
distribution may have either sign.  These local constraints do not guarantee
that the complete fitted polynomial is the log Laplace transform of a globally
nonnegative PDF.

\subsection{Gamma distribution as a diagnostic special case}
\label{app:gamma}

For comparison with classical path-distribution retrievals
\citep{funk2003,scholl2006}, suppose $p_0(\ell')$ is gamma distributed with mean
$m$ and shape $\kappa$:
\begin{equation}
  p_0(\ell')
  = \frac{(\kappa/m)^\kappa}{\Gamma(\kappa)}
    \ell'^{\kappa-1}
    \exp\left(-\frac{\kappa}{m}\ell'\right).
  \label{eq:app_gamma_pdf}
\end{equation}
Its Laplace transform is
\begin{equation}
  T(\tau)
  = \gamma\left(1+\frac{m}{\kappa}\tau\right)^{-\kappa}.
  \label{eq:app_gamma_transmittance}
\end{equation}
The ordinary cumulants are
$C_n=(n-1)!m^n/\kappa^{n-1}$, giving
\begin{equation}
  k_n=\frac{m^n}{\kappa^{n-1}},
  \qquad
  \kappa=\frac{k_1^2}{k_2},
  \qquad
  k_n=k_1\left(\frac{k_2}{k_1}\right)^{n-1}.
  \label{eq:app_gamma_mapping}
\end{equation}
The production estimator does not impose
Eq.~\eqref{eq:app_gamma_transmittance} or the coefficient sequence in
Eq.~\eqref{eq:app_gamma_mapping}; it fits $k_n$ independently.  Consequently,
$k_1^2/k_2$ is a gamma-equivalent shape diagnostic rather than evidence that
the inferred path distribution is gamma.  A gamma PDF cannot generally
represent a direct-beam-plus-scattered-tail mixture or distinct paths associated
with multiple cloud layers \citep{bennartz2006}.

\subsection{Layer-resolved generalization and interpretation limits}
\label{app:layered}

Without Eq.~\eqref{eq:app_collinear_paths}, the appropriate conservative-
ensemble expression is the multidimensional transform
\begin{equation}
  T_\lambda
  = \gamma_\lambda
    \int p_0(L_1,\ldots,L_Z)
    \exp\left[-\sum_z\alpha_{\lambda z}L_z\right]
    \mathrm{d}L_1\cdots\mathrm{d}L_Z.
  \label{eq:app_multidimensional}
\end{equation}
Different wavelengths sample different directions in the vector of
layer-dependent absorption coefficients.  The one-dimensional model can fail
when cloud-top reflection removes below-cloud paths, three-dimensional
transport adds path preferentially near cloud altitude, aerosols alter
lower-tropospheric paths, or line centers and wings have materially different
vertical absorption kernels.  Wavelength-dependent surface bidirectional
reflectance, polarization, and aerosol scattering can additionally violate the
common conservative-ensemble assumption.

These effects do not invalidate the trajectory-level equivalence theorem.
They limit the reduction from Eq.~\eqref{eq:app_general_equivalence} to
Eq.~\eqref{eq:app_laplace}.  We therefore use $k_1$ and $k_2$ as
band-effective absorber-path cumulants.  Exact identification with geometrical
path moments would require closure against layer-resolved Monte Carlo photon
histories.  Useful closure tests include predicting strong channels after
fitting weak and intermediate optical depths, comparing channels with similar
total $\tau$ but different vertical kernels, and fitting synthetic spectra for
which the conservative and layer-resolved path moments are known.
